\clearpage
\chapter{Results and Discussion}\label{chap:results_discussion}

This chapter presents the prediction results, matched representation comparisons,
computational costs and cue-level explanations. It follows BIU development with the
proposed XGBoost architecture across datasets, then neural integration, published-method
comparisons, cue explanations and efficiency. PIE provides the main empirical evidence.

\FloatBarrier
\section{Evaluation Scope}\label{sec:ch4_scope}

The dataset populations and splits follow Section~\ref{sec:method_datasets}; metric
and uncertainty calculations follow Section~\ref{sec:metrics_method}. PR-AUC denotes
average precision (AP). Tables report
three-run means unless stated otherwise; $\pm$ denotes sample standard deviation.
A displayed zero standard deviation can result from rounding.
Bold values mark the best displayed score within each classifier or dataset comparison,
including ties at the reported precision; lower error rates are preferred. Highlighting
is descriptive and does not indicate statistical significance or model selection. Appendix~\ref{app:experiment_map} links the descriptive model names to
experiment identifiers and source artifacts.

Confusion matrices show actual classes in rows and predicted classes in columns.
Their counts are averaged across runs and their percentages are calculated within each
actual class. Counts represent observation windows, so fractional averages do not
represent fractions of pedestrians.

The matched version comparisons revisit an already examined PIE test cohort and are
interpreted retrospectively. Bootstrap conclusions apply only to the comparisons for
which the intervals were computed; the reported rankings are not used for new model
selection.

% SOURCE: experiments/012_pie_biu_v2_30_60_benchmark/README.md
% SOURCE: experiments/096_pie_exp072_matched_biu_ablation/DESIGN.md

\FloatBarrier
\section{BIU Development and Representation Comparisons}\label{sec:ch4_initial}

\subsection{Initial BIU-v1 study: raw pose and full fusion}

The initial study tests whether BIUs improve prediction beyond BBox features, and
whether adding raw keypoints provides further benefit~\cite{Pasha2026BIU}.
Table~\ref{tab:ch4_pose_fusion} compares BBox alone, BBox+BIU-v1, BBox+raw keypoints
and their full fusion across four classifiers. Each row is an individual saved run.
This study predates BIU-v2; the following subsection evaluates the later direction-aware
representation. Raw keypoints include confidence values, unlike the XY-only neural
pose streams.
The recorded validation procedure selected Random Forest for all four input packages;
the other classifiers are shown to make their representation trade-offs visible.

% Source: experiments/aibthings_pie_val_fixed_0p5_selection/pie_val_fixed_0p5_all_models.csv
\begin{table}[!htbp]
\centering
\caption{Earlier PIE representation controls across four classifiers (individual runs).}\label{tab:ch4_pose_fusion}
\begin{tabularx}{\textwidth}{@{}Xrrr@{}}\toprule
\textbf{Input} & \textbf{\boldmath $F_1$} & \textbf{PR-AUC} & \textbf{ROC-AUC} \\[3pt]
\midrule
\thesistablegroup{4}{Logistic regression}
BBox & 0.811 & \textbf{0.808} & \textbf{0.939}\\[3pt]
BBox + BIU-v1 & \textbf{0.813} & 0.805 & \textbf{0.939}\\[3pt]
BBox + raw KP & 0.760 & 0.767 & 0.906\\[3pt]
BBox + raw KP + BIU-v1 & 0.760 & 0.759 & 0.909\\[3pt]
\midrule
\thesistablegroup{4}{Random Forest}
BBox & 0.865 & 0.886 & 0.961\\[3pt]
BBox + BIU-v1 & \textbf{0.869} & \textbf{0.912} & \textbf{0.965}\\[3pt]
BBox + raw KP & 0.812 & 0.856 & 0.952\\[3pt]
BBox + raw KP + BIU-v1 & 0.811 & 0.873 & 0.957\\[3pt]
\midrule
\thesistablegroup{4}{XGBoost}
BBox & \textbf{0.860} & 0.889 & 0.963\\[3pt]
BBox + BIU-v1 & 0.857 & \textbf{0.899} & \textbf{0.964}\\[3pt]
BBox + raw KP & 0.852 & 0.885 & 0.960\\[3pt]
BBox + raw KP + BIU-v1 & 0.853 & 0.883 & 0.963\\[3pt]
\midrule
\thesistablegroup{4}{LightGBM}
BBox & 0.856 & 0.876 & 0.960\\[3pt]
BBox + BIU-v1 & \textbf{0.860} & 0.892 & \textbf{0.963}\\[3pt]
BBox + raw KP & 0.857 & \textbf{0.893} & 0.961\\[3pt]
BBox + raw KP + BIU-v1 & 0.859 & 0.892 & 0.962\\[3pt]
\bottomrule
\end{tabularx}
\end{table}


For Random Forest, BIU+BBox attains $F_1=0.869$, compared with 0.865 for BBox alone and approximately 0.812
for either raw-pose fusion. Adding raw keypoints to BIU+BBox therefore does not recover an
advantage in this setting. Supplying more source coordinates is not sufficient to make
them useful to this learner. The result supports the usefulness of organizing pose into
the tested summaries, without implying that BIUs contain information absent from their
inputs or that raw pose is ineffective under every model.

For Random Forest, the BIU+BBox improvement over BBox is modest: the saved paired track-bootstrap interval
for their $F_1$ difference spans approximately $[-0.014,0.021]$. It does not establish a
statistically decisive gain. The more substantial descriptive difference from the
raw-pose arms should likewise not be assigned significance using an interval computed
for a different comparison.

LightGBM changes little across the four inputs: its $F_1$ ranges from 0.856 to 0.860.
Its highest PR-AUC occurs with raw keypoints, whereas its highest $F_1$ and ROC-AUC
occur with BIU-v1. For XGBoost, BIU-v1 improves the ranking metrics but slightly lowers
$F_1$ relative to BBox alone. These individual runs show why the Random Forest outcome
should not be generalized to every classifier.

% SOURCE: experiments/aibthings_pie_val_fixed_0p5_selection/pie_val_fixed_0p5_all_models.csv
% SOURCE: experiments/aibthings_pie_val_fixed_0p5_selection/pie_selected_by_val_fixed_0p5.csv
% SOURCE: experiments/aibthings_pie_val_fixed_0p5_selection/cluster_bootstrap_tests.csv

\FloatBarrier
\subsection{From BIU-v1 to BIU-v2}

The next comparison evaluates the direction-aware refinements introduced in BIU-v2
(Section~\ref{sec:biu_v2_method}). Table~\ref{tab:ch4_initial} compares BBox alone,
BBox+BIU-v1 and BBox+BIU-v2 within logistic regression, Random Forest and XGBoost.
These are three-seed means from the later version study, rather than the individual
runs in the preceding raw-pose comparison.

% Generated by experiments/thesis_chapter4_table_revision_2026_09_12/build_tables.py
\begin{table}[!htbp]
\centering

\caption{PIE classifier and BIU-version controls; $F_1$ includes sample SD.}\label{tab:ch4_initial}
\begin{tabularx}{\textwidth}{@{}Xrrr@{}}\toprule
\textbf{Predictor/input} & \textbf{\boldmath $F_1$} & \textbf{PR-AUC} & \textbf{ROC-AUC} \\[3pt]
\midrule
\thesistablegroup{4}{Logistic regression}
BBox & $0.811\pm0.000$ & 0.808 & 0.939\\[3pt]
BBox + BIU-v1 & $0.813\pm0.000$ & 0.805 & 0.939\\[3pt]
BBox + BIU-v2 & $\mathbf{0.816}\pm0.000$ & \textbf{0.818} & \textbf{0.940}\\[3pt]
\midrule
\thesistablegroup{4}{Random Forest}
BBox & $0.867\pm0.002$ & 0.887 & 0.961\\[3pt]
BBox + BIU-v1 & $\mathbf{0.868}\pm0.006$ & \textbf{0.911} & \textbf{0.965}\\[3pt]
BBox + BIU-v2 & $0.847\pm0.005$ & 0.900 & 0.962\\[3pt]
\midrule
\thesistablegroup{4}{XGBoost}
BBox & $0.849\pm0.002$ & 0.885 & 0.961\\[3pt]
BBox + BIU-v1 & $0.858\pm0.004$ & 0.904 & 0.964\\[3pt]
BBox + BIU-v2 & $\mathbf{0.860}\pm0.002$ & \textbf{0.916} & \textbf{0.967}\\[3pt]
\bottomrule
\end{tabularx}
\end{table}


For Random Forest, adding v1 increases mean PR-AUC from 0.887 to 0.911, while mean $F_1$
changes only from 0.867 to 0.868. The error counts explain the small operating-point gain:
false positives decrease from 186.0 to 161.0, but false negatives increase from 94.3 to
111.3. Substituting v2 does not improve this setting. Relative to v1, it produces
both more false positives and more false negatives, with mean $F_1$ falling to 0.847.

XGBoost gives a different result. Both BIU packages improve mean $F_1$ and ranking metrics
over its BBox reference. The v2 arm has slightly fewer FP and FN than v1 and the
highest PR-AUC within that classifier. Logistic regression also favors v2 on mean
$F_1$, although v1 produces one fewer false positive. Thus, v2 is a refinement whose
benefit depends on how the learner uses the representation. These comparisons do not
support a rule that the later version must replace v1 in every predictor.

Within XGBoost, the direction-aware BIU-v2 representation gives the highest mean
$F_1$, PR-AUC and ROC-AUC of the three input configurations. This supports its use
in the proposed multimodal architecture, while the Random Forest results retain the
case for v1 in that classifier. The following comparison examines the version choice with additional temporal and
contextual inputs.

Figure~\ref{fig:cm_initial} shows the XGBoost error trade-offs for these same three
input packages, complementing the classifier-wide metrics in Table~\ref{tab:ch4_initial}.
\begin{figure}[!htbp]
\centering
\includegraphics[width=\textwidth]{ch4_cm_initial.pdf}
\caption{Initial PIE XGBoost controls: mean confusion counts and within-class percentages.}\label{fig:cm_initial}
\end{figure}


\FloatBarrier
\section{Proposed Multimodal XGBoost Architecture}\label{sec:ch4_architecture}

\subsection{Matched BIU-version comparison in XGBoost}\label{sec:ch4_versions}

A further comparison examines the BIU versions with the temporal and ego-motion
inputs defined in Section~\ref{sec:lightweight_architecture}. It holds the PIE samples,
seeds and XGBoost recipe fixed while changing the representation package. Table~\ref{tab:matched_xgb_arms}
shows the shared baseline and the additions in each arm. The v2 package includes
its movement contrasts, so this compares complete feature packages rather than only
substituting one set of anatomical cues for another.

\begin{table}[!htbp]
\centering
\caption{Inputs to the matched XGBoost version comparison.}
\label{tab:matched_xgb_arms}


\begin{tabularx}{\textwidth}{@{}>{\RaggedRight\arraybackslash}p{0.12\textwidth}Xr@{}}
\toprule
\textbf{Arm} & \textbf{Contents} & \textbf{Dimensions} \\[3pt]
\midrule
Base & BBox53, six BBox contrasts, five ego summaries & 64 \\[3pt]
v1 & Base, 72 v1 features, three pose-quality contrasts & 139 \\[3pt]
v2 & Base, 56 v2 features, three pose-quality and twelve v2 contrasts & 135 \\[3pt]
\bottomrule
\end{tabularx}
\end{table}

% SOURCE: experiments/096_pie_exp072_matched_biu_ablation/scripts/feature_contract.py
% SOURCE: experiments/096_pie_exp072_matched_biu_ablation/DESIGN.md

Table~\ref{tab:ch4_versions} reports their performance. The no-BIU baseline is the
same baseline used in the PIE architecture results in Section~\ref{sec:ch4_xgb_pie}. False-positive and false-negative rates
show the error trade-offs: FPR is the percentage of noncrossing windows predicted as
crossing; FNR is the percentage of crossing windows missed.
% Generated by experiments/thesis_chapter4_table_revision_2026_09_12/build_tables.py
\begin{table}[!htbp]
\centering

\caption{Matched PIE BIU-version packages and their error trade-offs; $F_1$ includes sample SD.}\label{tab:ch4_versions}
\begin{tabularx}{\textwidth}{@{}Xrrrrr@{}}\toprule
\textbf{Input} & \textbf{\boldmath $F_1$} & \textbf{PR-AUC} & \textbf{ROC-AUC} & \textbf{FPR (\%)} & \textbf{FNR (\%)} \\[3pt]
\midrule
Baseline (64-D) & $\mathbf{0.877}\pm0.003$ & 0.919 & 0.970 & \textbf{5.81} & 9.45\\[3pt]
+ v1 package (139-D) & $0.873\pm0.003$ & 0.930 & 0.970 & 6.44 & 8.96\\[3pt]
+ v2 package (135-D) & $\mathbf{0.877}\pm0.002$ & \textbf{0.939} & \textbf{0.974} & 6.11 & \textbf{8.82}\\[3pt]
\bottomrule
\end{tabularx}
\end{table}


The v2 package has the highest mean PR-AUC and ROC-AUC. Relative to the base, its
seed-ensemble track-bootstrap intervals are positive for PR-AUC
($[0.0053,0.0369]$) and ROC-AUC ($[0.0008,0.0076]$), while the $F_1$ interval includes
zero ($[-0.0115,0.0184]$). These are intervals for paired ensemble comparisons; they are
not uncertainty intervals around the individual-seed means in the table.

Relative to v1, v2 has fewer mean false positives and false negatives. However, the
base still has a lower FPR (5.81\% versus 6.11\%), whereas v2 has a lower FNR
(8.82\% versus 9.45\%). A favorable ranking result
therefore does not establish fewer classification errors than a predictor without BIUs.
The initial classifier comparison also demonstrates classifier dependence.

BIU-v2 is therefore adopted as the principal behavioral representation of this thesis: its 56 features address
direction handling and explicit lower-body and torso motion with fewer inputs than v1's
72, and its package has the strongest ranking performance in the main matched PIE
comparison. This choice combines representation design, compactness and predictive
evidence.

\FloatBarrier

\subsection{PIE: behavioral model and road-context extension}\label{sec:ch4_xgb_pie}

Table~\ref{tab:ch4_architecture} starts with the no-BIU XGBoost baseline, then adds
the behavioral package and road context. The baseline contains 53 BBox summaries,
six BBox early/late features and five ego-speed summaries: $53+6+5=64$ inputs.
Adding BIU-v2 (56 features), 12 pose-derived movement contrasts and three pose-confidence
summaries gives 135 inputs. The 60-feature road branch brings the total to 195.
Section~\ref{sec:lightweight_architecture} describes these branches.

\begin{table}[!htbp]
\centering
\caption{PIE XGBoost: baseline, behavioral package and road extension (mean $\pm$ SD).}\label{tab:ch4_architecture}
\begin{tabularx}{\textwidth}{@{}Xrrr@{}}\toprule
\textbf{Predictor/input} & \textbf{\boldmath $F_1$} & \textbf{PR-AUC} & \textbf{ROC-AUC} \\[3pt]
\midrule
BBox + temporal + ego (\mbox{64-D}) & $0.877\pm0.003$ & $0.919\pm0.002$ & $0.970\pm0.001$\\[3pt]
Baseline + BIU-v2 package (\mbox{135-D}) & $0.877\pm0.002$ & $\mathbf{0.939}\pm0.001$ & $\mathbf{0.974}\pm0.000$\\[3pt]
Baseline + BIU-v2 package + road context (\mbox{195-D}) & $\mathbf{0.883}\pm0.004$ & $0.930\pm0.002$ & $0.973\pm0.001$\\[3pt]
\bottomrule
\end{tabularx}
\end{table}


The baseline and the XGBoost behavioral model both round to $F_1=0.877$, while the behavioral package
increases PR-AUC from 0.919 to 0.939 and ROC-AUC from 0.970 to 0.974.
With road context,
mean $F_1$ rises to 0.883, while PR-AUC decreases to 0.930 and ROC-AUC to 0.973.
The road branch therefore provides an operating-point trade-off, rather than an
improvement in every metric. Both configurations remain visible.
% SOURCE: experiments/096_pie_exp072_matched_biu_ablation/metrics/test_summary_results.csv
% SOURCE: experiments/116_xgboost_v2_core_completion/metrics/summary_results.csv

\FloatBarrier
\subsection{JAAD in-domain evidence}\label{sec:ch4_jaad}

Results are reported separately for clean event-aligned JAAD, API JAAD$_{beh}$ and
API JAAD$_{all}$, using the populations defined in Section~\ref{sec:method_datasets}.
Their event alignment and negative populations differ, so the comparisons below are
within each cohort.

\paragraph{In-domain prediction with BIU-v2.}
Table~\ref{tab:ch4_jaad_native} brings together the completed BBox controls and the
56-D BIU-v2 configurations. Every model is trained on its own JAAD training cohort.
The road-context inputs are BBox53, BIU-v2, 21 temporal contrasts, 60 compensated road
summaries and, where specified, 15 categorical vehicle-state summaries. The vehicle-only
clean control retains its recorded 145-input package. Rows are complete configurations;
no metric is taken from a different model to construct a best-score row.

\begin{table}[htbp]
\centering
\caption{JAAD in-domain XGBoost controls; $F_1$ includes sample SD. Dimensions count all predictor inputs.}\label{tab:ch4_jaad_native}
\begin{tabularx}{\textwidth}{@{}Xrrrr@{}}\toprule
\textbf{Inputs} & \textbf{Dim.} & \textbf{\boldmath $F_1$} & \textbf{PR-AUC} & \textbf{ROC-AUC} \\[3pt]
\midrule
\thesistablegroup{5}{Clean event-aligned JAAD}
BBox & 53 & $0.7038\pm0.0059$ & $0.6774$ & $0.7205$\\[3pt]
BBox + BIU-v2 & 109 & $0.7051\pm0.0008$ & $0.6743$ & $0.7080$\\[3pt]
BIU-v2 + BBox + temporal + vehicle & 145 & $0.7253\pm0.0020$ & $\mathbf{0.7304}$ & $\mathbf{0.7650}$\\[3pt]
BIU-v2 + BBox + temporal + road & 190 & $0.7191\pm0.0018$ & $0.6790$ & $0.7462$\\[3pt]
BIU-v2 + BBox + temporal + road + vehicle & 205 & $\mathbf{0.7282}\pm0.0088$ & $0.6806$ & $0.7527$\\[3pt]
\midrule
\thesistablegroup{5}{API JAAD$_{beh}$}
BBox & 53 & $0.7764\pm0.0026$ & $0.7225$ & $0.6512$\\[3pt]
BBox + BIU-v2 & 109 & $\mathbf{0.7782}\pm0.0044$ & $0.7460$ & $0.6885$\\[3pt]
BIU-v2 + BBox + temporal + road & 190 & $0.7714\pm0.0014$ & $0.7697$ & $0.6854$\\[3pt]
BIU-v2 + BBox + temporal + road + vehicle & 205 & $0.7755\pm0.0034$ & $\mathbf{0.7819}$ & $\mathbf{0.6970}$\\[3pt]
\midrule
\thesistablegroup{5}{API JAAD$_{all}$}
BBox & 53 & $0.6603\pm0.0040$ & $0.7495$ & $0.9278$\\[3pt]
BBox + BIU-v2 & 109 & $0.6772\pm0.0050$ & $0.7687$ & $0.9284$\\[3pt]
BIU-v2 + BBox + temporal + road & 190 & $0.7250\pm0.0038$ & $0.8112$ & $0.9453$\\[3pt]
BIU-v2 + BBox + temporal + road + vehicle & 205 & $\mathbf{0.7327}\pm0.0017$ & $\mathbf{0.8227}$ & $\mathbf{0.9487}$\\[3pt]
\bottomrule
\end{tabularx}
\end{table}


On clean JAAD, BBox+BIU-v2 changes mean $F_1$ only slightly, from 0.7038 to 0.7051,
while lowering AP and ROC-AUC. The road-context extension with vehicle state reaches
$F_1=0.7282$, AP 0.6806 and ROC-AUC 0.7527. The vehicle-only configuration has
lower $F_1$ (0.7253) but higher AP (0.7304) and ROC-AUC (0.7650). More context
therefore does not improve every metric, and a compact pose representation alone is
not enough to explain performance on this cohort.

On JAAD$_{beh}$, adding BIU-v2 to BBox improves all three mean metrics slightly.
The road-plus-vehicle model has higher AP (0.7819) and ROC-AUC (0.6970), but its
$F_1$ (0.7755) is below both BBox+BIU-v2 (0.7782) and BBox alone (0.7764).
On JAAD$_{all}$, the same road-context extension with vehicle state reaches $F_1=0.7327$, AP 0.8227
and ROC-AUC 0.9487, improving on the simpler BBox and BBox+BIU-v2 controls.
Its additional noncrossing population also changes the class balance: crossing windows
constitute 17.5\% of JAAD$_{all}$ test examples, compared with 62.6\% for JAAD$_{beh}$
and 48.0\% for clean JAAD. The all-cohort ranking result is consequently not evidence
that difficult behavior pedestrians have been resolved.

The 56-D representation is retained as the main definition, not declared a universal winner.
The capped PedCMT evidence in Section~\ref{sec:ch4_neural} supplies a separate
within-architecture comparison and is not pooled with these tree-model results.
% SOURCE: experiments/098_jaad_matched_biu_xgb/metrics/summary_results.csv
% SOURCE: experiments/100_jaad_exp099_biu_pathways/metrics/summary_results.csv
% SOURCE: experiments/102_jaad_api_matched_biu/metrics/summary_results.csv
% SOURCE: experiments/116_xgboost_v2_core_completion/metrics/summary_results.csv

JAAD supports conditional complementarity and
competitive performance on some cohort/metric combinations, alongside clear limitations
in behavior-cohort prediction. Zero-shot performance of the retained 56-D architecture
is not established by these in-domain experiments.


\FloatBarrier
\subsection{PSI: perceived-intention results}\label{sec:ch4_psi}

The PSI models use estimated ViTPose-S keypoints. Section~\ref{sec:method_psi} specifies the
protocol and the shared complete-input population used below. The classification task is inspired
by its dynamic human intent annotations~\cite{Jing2025PSI}. Fifteen observed frames
predict the aggregated intent label at the next frame. The test cohort contains 20,547 overlapping windows from
96 tracks: 15,462 positive and 5,085 negative windows. Windows advance by one frame,
so a sustained error on one track can contribute many erroneous windows. Raw error
counts are therefore much larger than counts of affected pedestrians and cannot be
compared directly with those from PIE's smaller test cohort. This experiment supplies
in-domain PSI evidence, not PIE-to-PSI zero-shot transfer.

% SOURCE: experiments/124_psi_matched_lightweight_ablation/metrics/summary_results.csv
\begin{table}[!htbp]
\centering
\caption{Matched PSI input ablation; $F_1$ is positive-class and includes sample SD.}\label{tab:ch4_psi}
\begin{tabularx}{\textwidth}{@{}Xrrrr@{}}\toprule
\textbf{Predictor/input} & \textbf{\boldmath $F_1$} & \textbf{\boldmath Macro-$F_1$} & \textbf{PR-AUC} & \textbf{ROC-AUC} \\[3pt]
\midrule
BBox53 & $0.746\pm0.005$ & $0.546$ & $0.833$ & $0.604$\\[3pt]
BBox53 + BIU-v2 (56-D) & $0.760\pm0.001$ & $0.570$ & $0.859$ & $0.652$\\[3pt]
XGBoost behavioral (135-D) & $0.804\pm0.001$ & $0.601$ & $0.847$ & $0.661$\\[3pt]
XGBoost + road context (195-D) & $\mathbf{0.825}\pm0.003$ & $\mathbf{0.630}$ & $\mathbf{0.862}$ & $\mathbf{0.690}$\\[3pt]
\bottomrule
\end{tabularx}
\end{table}


The matched comparison holds all training and evaluation windows fixed across the four
input configurations. Adding BIU-v2 to BBox raises mean macro-$F_1$ from 0.546 to
0.570, AP from 0.833 to 0.859 and ROC-AUC from 0.604 to 0.652. Adding temporal
and ego inputs together further raises macro-$F_1$ to 0.601, although AP falls to
0.847. Thus the sequence does not improve every metric at every step.

The road extension raises accuracy from 0.7043 to 0.7324, positive-class $F_1$ from
0.8040 to 0.8247, macro-$F_1$ from 0.6010 to 0.6297, AP from 0.8474 to 0.8622,
and ROC-AUC from 0.6608 to 0.6904. The behavioral model and road-context extension are reused unchanged;
the matched simple controls make the full 53/109/135/195-input sequence comparable
without changing the training population. The positive class accounts for about
75.3\% of test windows, so macro-$F_1$ and ranking metrics remain essential context.
% SOURCE: experiments/124_psi_matched_lightweight_ablation/metrics/summary_results.csv

\paragraph{Class-specific errors.}
Figure~\ref{fig:cm_psi} shows all four matched PSI arms on their shared test population.
The road extension reduces mean FNR from 19.4\% to 16.3\% and FPR from 60.5\%
to 58.4\%; substantial errors remain in both classes. Higher positive-class $F_1$ therefore does not imply equally
strong recognition of noncrossing intention. These counts refer to overlapping
windows with human-rated labels, not distinct pedestrians with verified future events.
\begin{figure}[!htbp]
\centering
\includegraphics[width=\textwidth]{ch4_cm_psi.pdf}
\caption{PSI perceived-intention models: mean confusion counts and within-class percentages.}\label{fig:cm_psi}
\end{figure}

\FloatBarrier

% SOURCE: experiments/114_psi_native_extended_xgb/test_evaluation/FINAL_REPORT.md
% SOURCE: experiments/113_psi_native_biu_xgb/test_evaluation/README.md
% SOURCE: experiments/113_psi_native_biu_xgb/test_evaluation/metrics/per_seed_results.csv
% SOURCE: experiments/113_psi_native_biu_xgb/DESIGN.md

\FloatBarrier
\subsection{Does ordered frame history improve the summaries?}\label{sec:ch4_ordered_history}

The development comparison in Section~\ref{sec:ordered_history_method} appends ordered
measurements from all 16 observed frames to the 135-feature model, producing 471 inputs.
It is evaluated on PIE validation data only. Across three seeds, mean $F_1$ changes from
0.7128 to 0.7110, PR-AUC from 0.7826 to 0.7831, and ROC-AUC from 0.9198 to 0.9186.
The small PR-AUC increase accompanies lower mean $F_1$ and ROC-AUC; there is no convincing
validation gain, and the candidate does not proceed to test evaluation.

This result supports retaining the compact summary input under the tested training recipe.
One possible explanation is that the summaries already capture much of the useful motion
within this short observation. The comparison does not establish that explanation:
extra measurements may be redundant or noisy, and the larger feature pool changes column
subsampling. It neither proves that summaries preserve all useful temporal information
nor rules out gains from another sequence representation or learner. These validation
scores are development evidence and are separate from the PIE test results in
Section~\ref{sec:ch4_xgb_pie}.

% SOURCE: experiments/118_pie_ordered_history_xgb/metrics/summary_results.csv
% SOURCE: experiments/118_pie_ordered_history_xgb/metrics/decision.json
% SOURCE: experiments/118_pie_ordered_history_xgb/metrics/INDEPENDENT_VERIFICATION.json

\FloatBarrier
\section{Complementarity in Neural Predictors}\label{sec:ch4_neural}

The neural evidence comprises the adapted MSLSTM representation controls and PedCMT with
a capped BIU correction. Sections~\ref{sec:mslstm_method} and~\ref{sec:pedcmt_method}
define their architectures and actual inputs, with the original methods providing research
lineage~\cite{Elsamalouty2025MSLSTM,Chen2024PedCMT}. Each comparison uses the thesis
anticipation benchmark and its own retained baseline. Initial PedCMT fusion trials provide
the design rationale in Chapter~\ref{chap:methodology}; the following results concern the
capped formulation.

\subsection{Adapted MSLSTM}

Table~\ref{tab:ch4_neural} groups the adapted MSLSTM comparisons on PIE and clean
JAAD. The PIE baseline follows Section~\ref{sec:mslstm_method}, with the clean-JAAD
adjustments in Section~\ref{sec:xgb_dataset_inputs}; the two
dataset groups use separately trained models.
% SOURCE: experiments/117_neural_v2_core_completion/metrics/summary_results.csv
% SOURCE: experiments/122_jaad_mslstm_no_noise/test_evaluation/metrics/summary_results.csv
\begin{table}[!htbp]
\centering

\caption{Adapted MSLSTM on PIE and clean JAAD; $F_1$ includes sample SD.}\label{tab:ch4_neural}
\begin{tabularx}{\textwidth}{@{}Xrrrr@{}}\toprule
\textbf{Predictor/input} & \textbf{Acc.} & \textbf{\boldmath $F_1$} & \textbf{PR-AUC} & \textbf{ROC-AUC} \\[3pt]
\midrule
\thesistablegroup{5}{PIE}
BBox53 + ego + YOLOP & $0.902$ & $0.810\pm0.009$ & $0.854$ & $0.935$\\[3pt]
Baseline + raw pose & $0.908$ & $0.823\pm0.009$ & $0.855$ & $0.937$\\[3pt]
Baseline + BIU-v1 (72-D) & $0.906$ & $0.822\pm0.003$ & $0.847$ & $0.935$\\[3pt]
Baseline + BIU-v2 (56-D) & $0.902$ & $0.812\pm0.006$ & $0.851$ & $0.932$\\[3pt]
Baseline + BIU-v2 + supporting cues (63-D) & $\mathbf{0.909}$ & $\mathbf{0.829}\pm0.005$ & $0.860$ & $0.938$\\[3pt]
\midrule
Baseline + GT Action/Look (annotation reference) & $0.908$ & $0.827\pm0.016$ & $\mathbf{0.863}$ & $\mathbf{0.946}$\\[3pt]
\midrule
\thesistablegroup{5}{Clean JAAD}
BBox53 + vehicle state + YOLOP & $0.713$ & $0.722\pm0.030$ & $0.796$ & $0.811$\\[3pt]
Baseline + raw pose & $0.724$ & $0.746\pm0.010$ & $0.788$ & $0.809$\\[3pt]
Baseline + BIU-v2 (56-D) & $0.712$ & $0.726\pm0.006$ & $\mathbf{0.800}$ & $0.813$\\[3pt]
Baseline + BIU-v2 + supporting cues (63-D) & $\mathbf{0.727}$ & $\mathbf{0.750}\pm0.008$ & $0.790$ & $\mathbf{0.815}$\\[3pt]
\bottomrule
\end{tabularx}
\end{table}


\paragraph{PIE.}
BIU-v2 (56-D) attains mean $F_1=0.812$, PR-AUC 0.851 and ROC-AUC 0.932.
The seven-feature neural extension improves these means to 0.829, 0.860 and 0.938,
respectively, with all three metrics higher in every paired seed. The extension therefore
helps this adapted recurrent predictor under the recorded training procedure.

The annotation reference adds ground-truth Action and Look to the same baseline.
It attains $F_1=0.827$, PR-AUC 0.863 and ROC-AUC 0.946. These are supplied
annotations, not outputs of deployed Action/Look detectors. BIU-v2 with supporting cues has slightly higher mean $F_1$, but lower ranking scores than this reference.

Both the 56-D and 63-D inputs remain in the table.
Raw pose produces fewer false positives than either the 56-D or 63-D BIU input,
so their advantages should not be interpreted as dominance over all controls and error types.
% SOURCE: experiments/117_neural_v2_core_completion/metrics/summary_results.csv
% SOURCE: experiments/117_neural_v2_core_completion/metrics/paired_deltas.csv


\paragraph{Clean JAAD.}
The clean-JAAD block in Table~\ref{tab:ch4_neural} presents the completed adaptation
with vehicle state and scene context. BIU-v2 plus supporting cues has the highest mean accuracy, $F_1$
and ROC-AUC, whereas BIU-v2 without supporting cues has the highest AP. The 63-D arm exceeds raw
pose on all four displayed means, but its AP remains below the no-pose baseline.
A higher mean therefore does not imply uniform improvement across metrics or seeds.

The initial configuration used noisy oversampling and selected checkpoints by validation
$F_1$. With that recipe, adding either BIU input reduced mean $F_1$, AP and ROC-AUC. The later comparison
uses the validation-AP and noise-free duplication recipe described in
Section~\ref{sec:mslstm_training}. The negative initial result remains evidence that
integration and training choices matter; the follow-up is not a causal ablation of all
changes between recipes. In particular, the seven supporting cues improve $F_1$ and
ROC-AUC over the 56-D follow-up arm while reducing AP.
% SOURCE: experiments/120_jaad_mslstm_v2_completion/FINAL_REPORT.md
% SOURCE: experiments/122_jaad_mslstm_no_noise/test_evaluation/metrics/summary_results.csv

Figure~\ref{fig:cm_mslstm} shows the PIE practical-input variants in Table~\ref{tab:ch4_neural}.
The ground-truth annotation reference is reported in the table only.
\begin{figure}[!htbp]
\centering
\includegraphics[width=\textwidth]{ch4_cm_mslstm.pdf}
\caption{Adapted MSLSTM on PIE: mean confusion counts and within-class percentages.}\label{fig:cm_mslstm}
\end{figure}


\FloatBarrier
\subsection{PedCMT with a capped BIU correction}

Table~\ref{tab:ch4_pedcmt_capped} evaluates the frozen-baseline correction defined in
Section~\ref{sec:pedcmt_capped}. It compares raw pose, BIU-v2 (56-D), and BIU-v2 plus
seven supporting cues (63-D), with the same parent and correction factor within each
dataset.

% SOURCE: experiments/119_pedcmt_v2_capped_completion/metrics/summary_results.csv
\begin{table}[!htbp]
\centering
\caption{Capped PedCMT with matched pose inputs; $F_1$ includes sample SD.}\label{tab:ch4_pedcmt_capped}
\begin{tabularx}{\textwidth}{@{}Xrrrr@{}}\toprule
\textbf{Predictor/input} & \textbf{Acc.} & \textbf{\boldmath $F_1$} & \textbf{PR-AUC} & \textbf{ROC-AUC} \\[3pt]
\midrule
\thesistablegroup{5}{PIE}
BBox + ego/vehicle baseline & $0.889$ & $0.793\pm0.009$ & $0.858$ & $0.919$\\[3pt]
Baseline + capped raw pose & $0.888$ & $0.792\pm0.010$ & $0.858$ & $0.920$\\[3pt]
Baseline + capped BIU-v2 (56-D) & $0.890$ & $0.795\pm0.009$ & $0.861$ & $0.922$\\[3pt]
Baseline + capped BIU-v2 + supporting cues (63-D) & $\mathbf{0.891}$ & $\mathbf{0.796}\pm0.011$ & $\mathbf{0.866}$ & $\mathbf{0.923}$\\[3pt]
\midrule
\thesistablegroup{5}{Clean JAAD}
BBox + ego/vehicle baseline & $0.633$ & $0.673\pm0.057$ & $0.626$ & $0.689$\\[3pt]
Baseline + capped raw pose & $0.639$ & $0.683\pm0.047$ & $0.602$ & $0.686$\\[3pt]
Baseline + capped BIU-v2 (56-D) & $0.648$ & $0.691\pm0.045$ & $0.657$ & $0.713$\\[3pt]
Baseline + capped BIU-v2 + supporting cues (63-D) & $\mathbf{0.654}$ & $\mathbf{0.698}\pm0.036$ & $\mathbf{0.670}$ & $\mathbf{0.726}$\\[3pt]
\bottomrule
\end{tabularx}
\end{table}


On PIE, the 63-D correction raises mean $F_1$ from 0.7926 to 0.7955, PR-AUC from
0.8584 to 0.8665, and ROC-AUC from 0.9192 to 0.9228. All three metrics improve over
both the parent and matched raw pose in every seed. The gains remain modest. The
56-D correction also improves these mean metrics, although its AP gain is not uniform
across seeds. Adding the supporting cues improves mean scores over the BIU-v2 alone.

On clean JAAD, the 63-D correction gives mean $F_1=0.6981$, PR-AUC 0.6700 and
ROC-AUC 0.7259, compared with 0.6734, 0.6256 and 0.6891 for the parent. Both BIU
forms exceed the raw-pose correction on all four displayed mean metrics. Their ranking
scores improve over the parent and raw pose in every seed; $F_1$ improves in two seeds
and ties in one. The 63-D input also has higher mean scores than the 56-D input.

These findings support the supporting-cue extension in the tested capped architecture,
not a universal need for all seven descriptors or an effect attributable to the cap
alone. No clustered confidence intervals were computed for these comparisons. The correction
factor transfers unchanged between datasets, but the parents and branches are trained
separately. Figure~\ref{fig:cm_pedcmt} shows their classification outcomes.
\begin{figure}[!htbp]
\centering
\includegraphics[width=\textwidth,height=0.85\textheight,keepaspectratio]{ch4_cm_pedcmt_capped.pdf}
\caption{Capped PedCMT on PIE and clean JAAD, with matched raw-pose, 56-D and 63-D corrections.}
\label{fig:cm_pedcmt}
\end{figure}

\FloatBarrier
% SOURCE: experiments/119_pedcmt_v2_capped_completion/metrics/summary_results.csv
% SOURCE: experiments/119_pedcmt_v2_capped_completion/metrics/paired_deltas.csv

\FloatBarrier
\section{Comparison with Published Methods}\label{sec:ch4_literature}

This section compares our results with published scores, grouped by dataset and cohort.
The TrajFusionNet rows retain the cited 2025 author-version scores; Chapter~\ref{chap:lit_rev}
cites the journal publication for the method.
These literature comparisons complement the controlled neural adaptations in
Section~\ref{sec:ch4_neural}; they are not matched reruns. Bold marks the highest
displayed value within each comparison, and a dash denotes an unreported metric.
AUC retains each source's definition; PSI's positive-class and macro-$F_1$ are reported
separately. Protocol differences are identified alongside each dataset comparison.

\subsection{PIE anticipatory benchmark}

Table~\ref{tab:ch4_literature_pie} compares the behavioral model and road-context extension of the proposed XGBoost architecture
with published approaches targeting 16 observations and TTE 30--60. PCPA supplies the
benchmark reference, GTransPDM supplies pose/without-pose variants, and TrajFusionNet
supplies a recent multimodal reference. Their rows come from their own tables; the latter
two use the cited author-manuscript versions. PCAP-SL is represented only by its selected
ViT-16-s configuration from Table 4 of the author-uploaded text. This is particularly
relevant because it also combines pretrained neural perception with XGBoost prediction.

\begin{table}[htbp]
\centering
\caption{Reported PIE results and our BIU-v2 XGBoost models (ours: three-run means).}\label{tab:ch4_literature_pie}
\begin{tabularx}{\textwidth}{@{}Xrrrrr@{}}\toprule
\textbf{Method} & \textbf{Acc.} & \textbf{\boldmath $F_1$} & \textbf{Precision} & \textbf{Recall} & \textbf{AUC} \\[3pt]
\midrule
PCPA~\cite{Kotseruba2021Benchmark} & 0.870 & 0.770 & -- & -- & 0.860\\[3pt]
GTransPDM: with keypoints~\cite{Xie2025GTransPDM} & 0.900 & 0.820 & 0.860 & -- & 0.870\\[3pt]
GTransPDM: without keypoints~\cite{Xie2025GTransPDM} & 0.920 & 0.860 & 0.850 & -- & 0.900\\[3pt]
TrajFusionNet~\cite{Landry2025TrajFusionNet} & 0.920 & 0.860 & 0.850 & 0.880 & 0.910\\[3pt]
PCAP-SL: selected ViT-16-s~\cite{Benkabou2026PCAPSL} & 0.934 & 0.863 & \textbf{0.898} & 0.830 & 0.896\\[3pt]
\midrule
Ours: XGBoost behavioral (\mbox{135-D}) & 0.9316 & 0.8770 & 0.8448 & 0.9118 & \textbf{0.9740}\\[3pt]
Ours: XGBoost + road context (\mbox{195-D}) & \textbf{0.9350} & \textbf{0.8829} & 0.8517 & \textbf{0.9164} & 0.9731\\[3pt]
\bottomrule
\end{tabularx}
\end{table}


The BIU-v2 road model attains mean $F_1=0.8829$ and accuracy 0.9350, compared with
0.863 and 0.934 for selected PCAP-SL. Its recall is higher (0.9164 versus 0.830),
while its precision is lower (0.8517 versus 0.898). Relative to TrajFusionNet's reported
$F_1=0.860$, the result is also numerically higher. These observations support competitive
performance with a different precision--recall balance, rather than superiority on every
measure. Within our models, the 135-D configuration has higher AP (0.9395 versus 0.9304)
and ROC-AUC (0.9740 versus 0.9731), while the road model has higher $F_1$; both complete
rows are retained.

\begingroup
\setlength{\emergencystretch}{2em}
Our results use three seeds and the PIE benchmark with separate routes. The GTransPDM
row also uses its benchmark split; its separate random-person evaluation is excluded. Exact PCAP-SL overlap,
generated sample identity and final-typeset table identity are unverified in the inspected
author copy. Threshold policies and AUC implementations are not established as identical
across papers; our AUC is computed from continuous probabilities. Accordingly, the
numerical differences are reported comparisons without a cross-paper significance claim.
The evidence for BIU contribution itself comes from the matched controls earlier in this
chapter. Prediction timing and upstream perception costs retain their separate measurement
boundaries in Section~\ref{sec:ch4_cost}.
\par\endgroup
% SOURCE: experiments/096_pie_exp072_matched_biu_ablation/metrics/test_summary_results.csv
% SOURCE: experiments/116_xgboost_v2_core_completion/metrics/summary_results.csv
% SOURCE: papers/Kotseruba_2021_PCPA_Benchmark.pdf, Table 3
% SOURCE: papers/Xie_2025_GTransPDM_arxiv_v2.pdf, Table I
% SOURCE: papers/Landry_Akhloufi_2025_TrajFusionNet_v1.pdf, Table II
% SOURCE: PCAP-SL author text, Table 4 (selected ViT-16-s)

\subsection{JAAD: comparison within each cohort}

Table~\ref{tab:ch4_literature_jaad} presents JAAD$_{beh}$ and JAAD$_{all}$ separately.
These API-based populations provide the closest supported comparison with published
16-frame, 1--2-second anticipation results. Clean JAAD uses our stricter event filtering;
its internal comparisons remain in Section~\ref{sec:ch4_jaad}.

% Reuses unchanged rows from ch4_literature_jaad_beh.tex and ch4_literature_jaad_all.tex.
\begin{table}[!htbp]
\centering
\caption{Published JAAD results by cohort; our scores are three-run means.}\label{tab:ch4_literature_jaad}
\begin{tabularx}{\textwidth}{@{}Xrrrrr@{}}\toprule
\textbf{Method} & \textbf{Acc.} & \textbf{\boldmath $F_1$} & \textbf{Precision} & \textbf{Recall} & \textbf{AUC} \\[3pt]
\midrule
\thesistablegroup{6}{JAAD$_{beh}$: behavior pedestrians}
PCPA~\cite{Kotseruba2021Benchmark} & 0.580 & 0.710 & -- & -- & 0.500\\[3pt]
Yang et al.~\cite{Yang2022FeatureFusion} & 0.620 & 0.740 & 0.650 & 0.850 & 0.540\\[3pt]
TrajFusionNet~\cite{Landry2025TrajFusionNet} & \textbf{0.740} & \textbf{0.820} & \textbf{0.730} & 0.930 & 0.680\\[3pt]
PCAP-SL: selected ViT-16-s~\cite{Benkabou2026PCAPSL} & \textbf{0.740} & 0.801 & 0.671 & \textbf{0.983} & \textbf{0.721}\\[3pt]
\midrule
Ours: BBox + BIU-v2 (109-D) & 0.6750 & 0.7782 & 0.6790 & 0.9114 & 0.6885\\[3pt]
Ours: XGBoost + road context + vehicle state (205-D) & 0.6729 & 0.7755 & 0.6795 & 0.9031 & 0.6970\\[3pt]
\midrule
\thesistablegroup{6}{JAAD$_{all}$: all pedestrians}
PCPA~\cite{Kotseruba2021Benchmark} & 0.850 & 0.680 & -- & -- & 0.860\\[3pt]
Yang et al.~\cite{Yang2022FeatureFusion} & 0.830 & 0.630 & 0.510 & 0.810 & 0.820\\[3pt]
TrajFusionNet~\cite{Landry2025TrajFusionNet} & 0.890 & 0.720 & 0.640 & 0.820 & 0.860\\[3pt]
PCAP-SL: selected ViT-16-s~\cite{Benkabou2026PCAPSL} & 0.852 & \textbf{0.841} & \textbf{0.791} & \textbf{0.892} & 0.852\\[3pt]
\midrule
Ours: BBox + BIU-v2 (109-D) & 0.8668 & 0.6772 & 0.5875 & 0.7992 & 0.9284\\[3pt]
Ours: XGBoost + road context + vehicle state (205-D) & \textbf{0.8918} & 0.7327 & 0.6447 & 0.8485 & \textbf{0.9487}\\[3pt]
\bottomrule
\end{tabularx}
\end{table}


\paragraph{Behavior pedestrians.}
Our BBox+BIU-v2 model reaches $F_1=0.7782$, above Yang et al.'s reported 0.740 but
below PCAP-SL's 0.801 and TrajFusionNet's 0.820. Adding road context and vehicle state
improves AP in our matched comparison but gives $F_1=0.7755$. Its precision is slightly
higher than PCAP-SL's (0.6795 versus 0.671), while recall is lower (0.9031 versus 0.983).
The added context therefore does not close the reported $F_1$ gap on this cohort.

\par\noindent\begin{minipage}{\textwidth}
\paragraph{All pedestrians.}
The road-context extension with vehicle state reaches accuracy 0.8918 and $F_1=0.7327$,
compared with TrajFusionNet's 0.890 and 0.720. PCAP-SL reports higher $F_1$ (0.841),
precision and recall, while our accuracy is higher than its 0.852. Performance is thus
competitive on some metrics, with a substantial $F_1$ gap to PCAP-SL.
\end{minipage}\par

The table uses PCPA and Yang et al.'s own tables, the cited TrajFusionNet author
manuscript, and PCAP-SL's selected ViT-16-s model from its author-uploaded Table 4.
The PCAP-SL source qualifications discussed for PIE also apply here. Exact generated
windows, decision rules and AUC implementations are not established as identical across
papers. In particular, our AUC uses continuous probabilities, so its higher JAAD$_{all}$
value is a reported numerical difference, not proof of a ranking advantage under an
identical evaluation.
% SOURCE: papers/Kotseruba_2021_PCPA_Benchmark.pdf, Table 3
% SOURCE: papers/Yang_2022_Feature_Fusion_Spatio_Temporal_Attention.pdf, Tables I--II
% SOURCE: papers/Landry_Akhloufi_2025_TrajFusionNet_v1.pdf, Table II
% SOURCE: PCAP-SL author full text, Table 4; URL and verification in thesis_jaad_finalization_2026_09_13/README.md

\subsection{PSI next-frame perceived intention}

Table~\ref{tab:ch4_literature_psi} uses the completed 56-D BIU-v2 models from
Section~\ref{sec:ch4_psi}, including the matched 109-, 135- and 195-input configurations.
It compares them with the published PSI results~\cite{Jing2025PSI}.
Both evaluations use 15 observations and
predict intent at frame 16, which is separate from the PIE/JAAD future-crossing-event
benchmark. The table keeps macro-$F_1$, positive-class $F_1$ and class-wise average
accuracy distinct.

\begin{table}[htbp]
\centering
\caption{PSI results under each study's evaluation protocol; BA: balanced accuracy.}\label{tab:ch4_literature_psi}
\begin{tabularx}{\textwidth}{@{}Xrrrr@{}}\toprule
\textbf{Method} & \textbf{Acc.} & \textbf{\boldmath $F_1^{+}$} & \textbf{\boldmath Macro-$F_1$} & \textbf{BA} \\[3pt]
\midrule

eP2P~\cite{Jing2025PSI} & 0.6030 & -- & 0.5701 & 0.6226\\[3pt]
eP2P + disagreement + reasoning~\cite{Jing2025PSI} & 0.6309 & -- & 0.5855 & 0.6253\\[3pt]
\midrule
Ours: BBox + BIU-v2 (109-D) & 0.6538 & 0.7600 & 0.5696 & 0.5777\\[3pt]
Ours: XGBoost behavioral (135-D) & 0.7043 & 0.8040 & 0.6010 & 0.6005\\[3pt]
Ours: XGBoost + road context (195-D) & \textbf{0.7324} & \textbf{0.8247} & \textbf{0.6297} & \textbf{0.6261}\\[3pt]
\bottomrule
\end{tabularx}
\end{table}


The road-context extension has accuracy 0.7324, macro-$F_1=0.6297$ and balanced accuracy
0.6261. The PSI eP2P model with disagreement and reasoning reports 0.6309, 0.5855
and 0.6253, respectively. Our accuracy and macro-$F_1$ are numerically higher, while
balanced accuracy is nearly the same.

Sample construction limits the interpretation of these differences. The PSI paper
uses a 75/5/20 video split and 60-frame clips with overlap 0.9. Our recovered release lists
contain 108/35/49 train/validation/test videos and use 16-frame windows at stride one,
with 15 observed frames and the target at frame 16. The shared population for our matched configurations is defined in
Section~\ref{sec:method_psi}. The table therefore supplies contextual
evidence, not an identical-benchmark
leaderboard.
The supported PSI finding is improved discrimination across both intent classes under
the matched representation/context controls, with substantial remaining errors.
% SOURCE: NeurIPS 2025 PSI paper, Table 4 and Section 3.2 (p. 9).
% https://papers.neurips.cc/paper_files/paper/2025/file/436fb0fa57c75e0d2063b5bc19a21da1-Paper-Datasets_and_Benchmarks_Track.pdf
% SOURCE: experiments/116_xgboost_v2_core_completion/metrics/summary_results.csv
% SOURCE: experiments/113_psi_native_biu_xgb/DESIGN.md
% SOURCE: experiments/114_psi_native_extended_xgb/APPROVED_COHORT_AMENDMENT.md
% SOURCE: experiments/124_psi_matched_lightweight_ablation/DESIGN.md

% SOURCE: experiments/thesis_confusion_matrices_2026_09_12/sources_sha256.json
% SOURCE: experiments/thesis_confusion_matrices_2026_09_12/matrix_cells.csv

\FloatBarrier
\section{Interpreting BIU Importance and Cue Contributions}\label{sec:ch4_inspectability}

This section asks two questions: how much the fitted predictor relies on BIUs, and
how often each cue pushes its score toward the known label. Importance describes
magnitude; toward-label frequency describes direction. Neither is the percentage of
correct predictions caused by BIUs. Section~\ref{sec:cue_attribution_method} defines
the calculations. The plots explain retained models without retraining or selecting
features from their test explanations.

\subsection{BIU-v2 reliance across datasets}

Figure~\ref{fig:ch4_v2_input_share} reports the 56-feature BIU-v2 block's share of
summed absolute feature SHAP contributions. The additional early/late pose contrasts remain
outside this block. The bars refer to the 135-input PIE behavioral model, the
145-input clean-JAAD vehicle model and the 135-input PSI behavioral model.
All use the same attribution calculation, but their populations, available context
and, for PSI, prediction task differ.

\begin{figure}[!htbp]
\centering
\includegraphics[width=\textwidth]{ch4_v2_input_share.pdf}
\caption{BIU-v2 share of absolute feature attribution on test data; three-seed means.}
\label{fig:ch4_v2_input_share}
\end{figure}

The share is 25.6\% on PIE, 23.1\% on clean JAAD with vehicle context and 31.4\%
on PSI. BIUs therefore account for a substantial part of the explanations in all
three settings. A higher share means greater attribution to this block within that
predictor; it does not establish a better representation or higher prediction accuracy.

\subsection{Cue contributions across datasets}

Overall prediction scores tell us how well the model performs, but they do not show
which observed behaviors influence its decisions. Examining individual BIUs makes
this evidence visible. Importance identifies the cues with the strongest influence,
while toward-label frequency shows how consistently their contributions support the
known outcome. Together, they help direct error analysis toward influential cues
whose contributions frequently oppose the label.

Table~\ref{tab:ch4_v2_cue_importance} shows the share of within-BIU absolute
attribution assigned to each cue. Figure~\ref{fig:ch4_v2_cue_direction} shows every v2 cue's toward-label frequency
for these three settings. It answers how often a cue supports the annotation,
including windows whose final prediction is still incorrect.

\begin{table}[!htbp]
\centering
\caption{Within-BIU absolute attribution by cue (\%); three-seed means.}
\label{tab:ch4_v2_cue_importance}
\begin{tabularx}{\textwidth}{@{}Xrrr@{}}
\toprule
\textbf{Cue} & \textbf{PIE} & \shortstack{\textbf{JAAD clean}\\\textbf{+ vehicle}} & \textbf{PSI}\\
\midrule
Torso lean & 3.8 & 9.9 & 3.1\\
Head orientation & 9.7 & 11.7 & \textbf{38.1}\\
Shoulder-line orientation & 7.2 & 9.2 & 6.5\\
Stride length & 10.0 & 12.0 & 20.5\\
Ankle elevation & \textbf{17.9} & \textbf{19.2} & 1.6\\
Hip-center motion & 2.0 & 4.9 & 2.1\\
Knee-center motion & 6.4 & 3.3 & 3.0\\
Ankle-center motion & 2.2 & 3.3 & 5.9\\
Lower-body center motion & 8.0 & 1.1 & 6.3\\
Forward ankle displacement & 2.8 & 8.0 & 1.4\\
Relative ankle displacement & 6.7 & 2.5 & 6.2\\
Lower-body motion magnitude & 1.1 & 2.2 & 2.5\\
Torso-center motion & 5.2 & 3.5 & 0.5\\
Shoulder-center motion & 16.4 & 3.5 & 0.8\\
Torso motion magnitude & 0.8 & 5.4 & 1.4\\
\bottomrule
\end{tabularx}
\end{table}


\begin{figure}[!htbp]
\centering
\includegraphics[width=\textwidth]{ch4_v2_cue_direction.pdf}
\caption{Percentage of test windows in which each BIU-v2 cue supports the label; three-seed means.}
\label{fig:ch4_v2_cue_direction}
\end{figure}

Ankle elevation has the largest within-BIU importance on clean JAAD with vehicle
context (19.2\%), yet it points toward the label in only 47.3\% of windows.
Its magnitude and supportive frequency therefore give different descriptions of its
use. On PIE, shoulder-center motion has
16.4\% of within-BIU importance and supports the label in 86.2\% of windows.
These observations show why a single importance ranking is insufficient for deciding
whether a cue is helpful or misleading.

% SOURCE: experiments/matched_biu_cue_rankings_2026_09_12/summary.csv
% SOURCE: experiments/matched_biu_cue_rankings_2026_09_12/DESIGN.md
% SOURCE: experiments/matched_biu_cue_rankings_2026_09_12/verification.json
% SOURCE: experiments/aibthings_pie_val_fixed_0p5_selection/rf_biu_bbox_group_importances.csv

On PSI, the retained behavioral model in Table~\ref{tab:ch4_psi} is explained using
all 20,547 test windows and three seeds. Head orientation accounts for 38.1\%
of the within-BIU absolute attribution, followed by stride length at 20.5\%.
These cues support the label in 65.2\% and 59.8\% of windows, respectively.

Class-specific inspection also matters: head orientation supports the label in
69.9\% of positive-label windows and 50.9\% of negative-label windows. The overall
percentage reflects PSI's class balance as well as the model's behavior. These
contributions describe support for perceived-intention labels, not verified future
crossing events.
% SOURCE: experiments/psi_v2_attribution_2026_10_19/summary.csv
% SOURCE: experiments/psi_v2_attribution_2026_10_19/verification.json

For adaptation to another dataset, these comparisons identify cues worth closer
inspection: for example, a cue with strong label support in one setting but weaker
support in another. Its pose measurements and associated prediction errors can then
be examined to understand the difference. This is a practical benefit of named BIUs:
model analysis can refer to specific behaviors such as shoulder motion or ankle
elevation. Collection conditions, pose quality, class composition, available context
and the prediction task remain relevant when interpreting these differences.

\FloatBarrier
\section{Measured Downstream Cost}\label{sec:ch4_cost}

Table~\ref{tab:ch4_runtime} reports the measured BIU-v2 pipeline cost on an Intel Core
i7-11800H CPU. From ready features, prediction averages 1.06\,ms for the behavioral model
and 1.17\,ms for the road-context extension. Recomputing the behavioral model’s features from available boxes,
keypoints and speed, then predicting, takes 9.55\,ms per observation window
(Figure~\ref{fig:ch4_runtime}). The corresponding 95th percentile is 10.59\,ms.

\begin{table}[tbp]
\centering
\caption{Measured CPU milliseconds per observation window.}\label{tab:ch4_runtime}
\begin{tabularx}{\textwidth}{@{}Xrrr@{}}\toprule
\textbf{Model} & \shortstack{\textbf{Prediction}\\\textbf{only}\\\textbf{Mean (ms)}} & \shortstack{\textbf{Features +}\\\textbf{prediction}\\\textbf{Mean (ms)}} & \shortstack{\textbf{Features +}\\\textbf{prediction}\\\textbf{95th pct. (ms)}}\\[3pt]\midrule
XGBoost behavioral (135-D) & 1.06 & 9.55 & 10.59\\[3pt]
XGBoost + road context (195-D) & 1.17 & 9.64 & 10.33\\[3pt]
\bottomrule
\end{tabularx}
\smallskip
\begin{minipage}{\textwidth}\footnotesize
Ready-feature means average three model seeds; sequential feature-to-prediction timings
use seed 42 and include BBox extraction. The road row starts from available road geometry.
\end{minipage}
\end{table}

\begin{figure}[tbp]
\centering\begingroup\setstretch{1.0}
\begin{tikzpicture}[x=1cm,y=1cm,>=Latex,
 every node/.style={font=\small,align=center},
 b/.style={draw=black!55,fill=black!3,text width=2.55cm,minimum height=0.95cm},
 f/.style={draw=blue!55!black,fill=blue!5,text width=3.05cm,minimum height=0.95cm},
 w/.style={draw=blue!55!black,fill=blue!5,text width=6.9cm,minimum height=0.95cm}]

\fill[blue!5] (0.7,-3.5) rectangle (9.7,0.5);
\node[font=\Large\bfseries,text=blue!55!black] at (5.2,-0.1) {9.55 ms / observation};
\node[font=\small] at (5.2,-0.7) {Measured behavioral pipeline on CPU};
\node[b,text width=1.7cm] (box) at (-0.8,-1.7) {Boxes + ego speed};
\node[b,text width=1.7cm] (pose) at (-0.8,-3) {Keypoints};
\node[text width=2.1cm] (biu) at (2.5,-2.9) {\textbf{Feature}\\construction};
\node[text width=2.1cm] (fusion) at (5.3,-2.9) {\textbf{Input}\\packing};
\node[text width=2.1cm] (rf) at (8.2,-2.9) {\textbf{XGBoost}\\prediction};
\draw[->] (pose) -- (biu); \draw[->] (biu) -- (fusion); \draw[->] (fusion) -- (rf);
\draw[->] (box.east) -- (2.5,-1.7) -- (biu.north);
\node[font=\footnotesize] at (4,-4) {Intel Core i7-11800H CPU; one observation per prediction};

\end{tikzpicture}\endgroup
\caption{Measured BIU-v2 XGBoost behavioral pipeline from available boxes, keypoints and ego speed.}\label{fig:ch4_runtime}
\end{figure}


Most execution time is spent constructing features. BIU-v2 extraction averages 2.69\,ms,
and the early/late contrasts, which rerun the extractor on both halves, average 3.98\,ms.
The road-context extension takes 9.64\,ms when road-relative geometry is already
available; deriving that geometry from masks is not included. These measurements support
a lightweight downstream decision stage, without assigning its latency to upstream
neural perception.

% SOURCE: experiments/v2_xgboost_runtime_2026_09_13/headline.csv
% SOURCE: experiments/v2_xgboost_runtime_2026_09_13/stage_summary.csv

\paragraph{Competitive performance with a compact predictor.}
Table~\ref{tab:ch4_model_size} compares the retained predictors under the separate
ready-input CPU audit defined in Section~\ref{sec:efficiency_method}. Both XGBoost
configurations contain 250 trees: 5,674--5,756 total nodes for the behavioral model and
5,872--5,890 for the road-context extension across seeds. The road-context extension and its imputer occupy
0.360--0.361\,MiB and take a mean 1.174\,ms from ready inputs. This is distinct from
the 9.64\,ms feature-construction-plus-prediction measurement above.

Within adapted MSLSTM, adding the 63-D BIU input increases the baseline by 74,241
parameters, including its fusion weight, compared with 2,240,770 for the raw-pose pathway.
The complete MSLSTM predictor with BIUs is also faster than its raw-pose counterpart
in this ready-input audit, although
BIU construction is excluded. The 56-D and 63-D projectors differ by only 896 weights.
The completed capped PedCMT has 1,834,416 frozen parent parameters and either 90,113
(56-D) or 91,009 (63-D) trainable correction parameters. Reporting only the correction
would omit most of the deployed predictor. Its newer timing run uses prepared validation
windows on the same CPU protocol; it is distinct from the earlier test-window audit.
Inference-only file sizes for these new branches were not measured in that audit.

These measurements support a compact and fast downstream XGBoost predictor among the
retained implementations. Published-model performance remains a separate comparison
in Section~\ref{sec:ch4_literature}.

\begin{table}[!htbp]
\centering
\caption{Retained predictor complexity and ready-input CPU latency.}\label{tab:ch4_model_size}
\small
\begin{tabularx}{\textwidth}{@{}Xrrrr@{}}\toprule
\textbf{Predictor} & \shortstack{\textbf{Trees /}\\\textbf{parameters}} & \shortstack{\textbf{Size}\\\textbf{(MiB)}} & \shortstack{\textbf{Mean}\\\textbf{(ms)}} & \shortstack{\textbf{p95}\\\textbf{(ms)}}\\[3pt]
\midrule
XGBoost behavioral (\mbox{135-D}) & 250 trees & 0.351--0.353 & 1.065 & 1.170\\[3pt]
XGBoost + road context (\mbox{195-D}) & 250 trees & 0.360--0.361 & 1.174 & 1.258\\[3pt]
MSLSTM baseline & 6,765,321 & 25.829 & 7.749 & 8.302\\[3pt]
MSLSTM + BIU-v2 (56-D) & 6,838,666 & 26.110 & 7.780 & 8.322\\[3pt]
MSLSTM + BIU-v2 + supporting cues (63-D) & 6,839,562 & 26.114 & 7.965 & 8.648\\[3pt]
MSLSTM + raw pose & 9,006,091 & 34.383 & 10.443 & 11.094\\[3pt]
Capped PedCMT + BIU-v2 (56-D) & 1,924,529 & -- & 3.424 & 3.770\\[3pt]
Capped PedCMT + BIU-v2 + supporting cues (63-D) & 1,925,425 & -- & 3.430 & 3.789\\[3pt]
\bottomrule
\end{tabularx}
\smallskip
\begin{minipage}{\textwidth}\footnotesize
Tree counts and neural parameter counts use different units. Sizes describe uncompressed
files, not RAM; PedCMT inference-only sizes were not measured (--). Complete predictors
are timed, including the frozen PedCMT parent. Section~\ref{sec:efficiency_method}
defines the shared CPU protocol and the different sampled windows.
\end{minipage}
\end{table}

% SOURCE: docs/retained_predictor_complexity_2026_09_16/evidence/AUDIT.md
% SOURCE: docs/retained_predictor_complexity_2026_09_16/evidence/latency_summary.csv
% SOURCE: docs/retained_predictor_complexity_2026_09_16/evidence/serialized_sizes.csv

\FloatBarrier
\section{Limitations of the Current Evidence}\label{sec:ch4_limits}

The gains vary with the classifier, available inputs and dataset. BIU-v2 gives strong
results in the proposed architecture, but no representation improves every reported
metric. The comparisons therefore show where each configuration is useful.

BIUs depend on estimated two-dimensional keypoints. Occlusion, viewpoint and pose
errors can affect their measurements. Named cues and SHAP make the predictor's use
of these measurements easier to inspect; they do not explain what caused the
pedestrian's behavior.

The main comparisons evaluate models within each dataset; they do not establish
performance in new environments. PSI uses perceived-intention labels rather
than the PIE/JAAD future-crossing task. Overlapping windows from one pedestrian are
related observations; the selected track-level analyses account for this dependence.
Published-method comparisons remain qualified by differences in sample construction
and evaluation settings.
\clearpage
